% Created 2022-02-05 sáb 18:16
% Intended LaTeX compiler: pdflatex
\documentclass[11pt]{article}
\usepackage[utf8]{inputenc}
\usepackage[T1]{fontenc}
\usepackage{graphicx}
\usepackage{longtable}
\usepackage{wrapfig}
\usepackage{rotating}
\usepackage[normalem]{ulem}
\usepackage{amsmath}
\usepackage{amssymb}
\usepackage{capt-of}
\usepackage{hyperref}
\usepackage{minted}
\author{Laura Brustenga i Moncusí}
\date{\today}
\title{}
\hypersetup{
 pdfauthor={Laura Brustenga i Moncusí},
 pdftitle={},
 pdfkeywords={},
 pdfsubject={},
 pdfcreator={Emacs 28.0.90 (Org mode 9.6)}, 
 pdflang={English}}
\begin{document}

\tableofcontents

\#-\textbf{- mode: org -}-
96120: THE DEGREE OF THE LINEAR ORBIT OF A CUBIC SURFACE
(join work with SASCHA TIMME and MADELEINE WEINSTEIN)



\section{Setting}
\label{sec:org18aef96}
\subsection{Base field: \(\mathbb{C}\)}
\label{sec:orgf5b49a2}
\subsection{Cubic surfaces \(S\subseteq\mathbb{P}^3_{x:y:z:w}\)}
\label{sec:orga0e684a}

\(S=V_+(F)\subseteq \mathbb{P}^3\) with \(F=a_{300}x^3+a_{210}x^2y+\dots\)

There are 20 coefficients, so its parameter space is \(F=[a_{300}:\ldots:a_{003}]\in \mathbb{P}^{19}\)

\subsection{Linear change of coordinates \(\varphi\in\mathrm{PGL}(\mathbb{C},4)\) acting on \(\mathbb{P}^{19}\)}
\label{sec:org897f374}

\(\varphi : \mathbb{P}^3\to \mathbb{P}^3'\) defines an action on \(\mathbb{P}^{19}\)

$$
\varphi = \left( \begin{array}{ccccc}
A_{1,1} & A_{1,2} & A_{1,3} & A_{1,4} \\
A_{2,1} & A_{2,2} & A_{2,3} & A_{2,4} \\
A_{3,1} & A_{3,2} & A_{3,3} & A_{3,4} \\
A_{4,1} & A_{4,2} & A_{4,3} & A_{4,4} \\
\end{array}\right)\in\mathrm{PGL}(\mathbb{C},4)\subseteq \mathbb{P}^{4\times 4 -1}
$$

$$\tilde{\varphi}: \mathbb{C}[x',y',z',w']\to \mathbb{C}[x,y,z,w]$$

sending \(x'\) to \(A_{1,1}x+A_{1,2}y+A_{1,3}z+A_{1,4}w\)
and so on.

We have \(\varphi(S) = V_+(F\circ \tilde{\varphi})\) is another cubic surface.

So, \(\mathrm{PGL}(\mathbb{C},4)\) acts on \(\mathbb{P}^{19}\) by precomposition

\(\varphi\cdot F := F\circ \tilde{\varphi}\)
\subsubsection{Example 2.1}
\label{sec:orgeb9f6ee}

Pairs of points in \(\mathbb{P}^1_{x:y}\)

$$\tilde{\varphi} : \mathbb{C}[x,y]\to \mathbb{C}[x,y]$$

$$ \varphi\cdot f = \left( \begin{array}{ccccc}
                             a_{11}^2 & a_{11}a_{21} & a_{21}^2 \\
                             2a_{11}a_{12} & a_{11}a_{22}+a_{12}a_{21} & 2a_{21}a_{22} \\
                             a_{12}^2 & a_{12}a_{22} & a_{22}^2 \\
                           \end{array}\right) \left(
                           \begin{array}{c}
                             b_1\\
                             b_2\\
                             b_3\\
\end{array}\right)$$

\subsubsection{More compactly}
\label{sec:org960437c}

We consider a \(\mathbb{C}\) vector space \(V\) of dimension 4.
Given
\(\varphi\in \mathrm{PGL}(V)\) and \(F\in \mathbb{P} \mathrm{Sym}^3(V)\),

$$\varphi\cdot  F := \mathrm{Sym}^3(\varphi)(F)$$

\subsection{The variety \(\Omega_{F}\), closure of \(\mathrm{PGL}(\mathbb{C},4)\cdot F\)}
\label{sec:org8c4aa2e}

We do \emph{not} know its equations, but we have the morphism

$$ \Theta_F: \mathrm{PGL}(\mathbb{C},4) \to \Omega_F \subseteq \mathbb{P}^{19}$$

sending \(\varphi\) to \(\varphi\cdot F\)

\begin{itemize}
\item It is \emph{dominant} (trivial)

\item It is one-to-one (a general \(F\) has trivial stabilizer)

\begin{itemize}
\item Let me call it ``the parametrization''

\item \(\dim\big(\Omega_{F}\big) = \dim \big(\mathrm{PGL}(\mathbb{C},4)\big) = 15\)
\end{itemize}
\end{itemize}
\subsection{Goal: Compute the degree of \(\Omega_{F}\)}
\label{sec:orga595550}
Motivation: revisit cubic surfaces with modern techniques.
\begin{itemize}
\item challenge modern techniques
\item what can be said about
\end{itemize}

Question 1 of

\href{https://lematematiche.dmi.unict.it/index.php/lematematiche/issue/view/82}{TWENTY-SEVEN QUESTIONS ABOUT THE CUBIC SURFACE (KRISTIAN RANESTAD - BERND STURMFELS)}

\section{Formulations}
\label{sec:orgef81a00}
\subsection{Common part}
\label{sec:orgf6a1ed9}
\subsubsection{Packages}
\label{sec:orgf95c055}

\begin{minted}[breaklines,linenos,numbersep=3pt,mathescape=true,frame=lines,framesep=2mm]{julia}
using HomotopyContinuation
using LinearAlgebra, PolynomialRoots, Combinatorics
\end{minted}

\subsubsection{Coordinates of \(\mathbb{P}^3\) and \(\mathrm{PGL}(\mathbb{C},4)\)}
\label{sec:org904b856}

\begin{minted}[breaklines,linenos,numbersep=3pt,mathescape=true,frame=lines,framesep=2mm]{julia}
@var x y z w
@var A[1:4,1:4]
\end{minted}

\subsubsection{General cubic surface}
\label{sec:org8ad6385}

\begin{minted}[breaklines,linenos,numbersep=3pt,mathescape=true,frame=lines,framesep=2mm]{julia}
monomials = [prod([x,y,z,w].^exp) for exp in multiexponents(4, 3)]
c₀ = randn(ComplexF64, 20);
F = sum(c₀ .* monomials);
\end{minted}

\subsubsection{The ``parametrization'' \(\Theta_F: \mathrm{PGL}(\mathbb{C},4) \to \Omega_F\)}
\label{sec:orga2a635f}

\begin{minted}[breaklines,linenos,numbersep=3pt,mathescape=true,frame=lines,framesep=2mm]{julia}
Θ = coefficients(subs(F, [x,y,z,w] => A*[x,y,z,w]), [x,y,z,w])
typeof(Θ)
size(Θ)
\end{minted}

\subsection{Formulation 1: 15 general hyperplanes}
\label{sec:org75d8832}
\subsubsection{General linear space in \(\mathbb{P}^{19}\) of dim 4}
\label{sec:orgdfd9997}

\begin{minted}[breaklines,linenos,numbersep=3pt,mathescape=true,frame=lines,framesep=2mm]{julia}
L = randn(ComplexF64, 15, 20);
\end{minted}

\subsubsection{Pulling back to \(\mathrm{PGL}(\mathbb{C},4)\)}
\label{sec:org03be674}

\begin{minted}[breaklines,linenos,numbersep=3pt,mathescape=true,frame=lines,framesep=2mm]{julia}
Eq1 = L*Θ;
size(Eq1)
\end{minted}

\subsubsection{Solving}
\label{sec:org4f132b6}

\begin{minted}[breaklines,linenos,numbersep=3pt,mathescape=true,frame=lines,framesep=2mm]{julia}
System1 = System(Eq1; variables=vec(A));
\end{minted}

Bezout bound
\begin{minted}[breaklines,linenos,numbersep=3pt,mathescape=true,frame=lines,framesep=2mm]{julia}
solve(System1,start_system=:total_degree)
\end{minted}

Tracking 14348907 paths\ldots{}   0\%|                        |  ETA: 29.85 days

BKK bound
\begin{minted}[breaklines,linenos,numbersep=3pt,mathescape=true,frame=lines,framesep=2mm]{julia}
solve(System1,start_system=:polyhedral)
\end{minted}

ERROR: InterruptException:

\subsubsection{Pros and Cons}
\label{sec:org212474c}

\begin{itemize}
\item (con) In fact, the build \(\Theta\) has source \(\mathbb{P}^{15}\)
It is a rational morphims (so, System1 contain its base locus: \(\det(A)=0\))

\begin{itemize}
\item Classify solutions. Problem: pick tolerance
\item Use monodromy method. Problem: which parameters
\end{itemize}

\item (pro) Better suited for the trace test (we will see it later).
\end{itemize}

\subsection{Formulation 2: 15 general points (monodromy)}
\label{sec:orgadba504}

\subsubsection{Idea}
\label{sec:orge570c2a}

Given a point \(p=[x_0:y_0:z_0:w_0]\in \mathbb{P}^3\),
the set of cubics passing through \(p\) form a hyperplane in \(\mathbb{P}^{19}\)

\(\Pi_p=V_{+}(x_0^3a_{3000}+x_0^2y_0a_{2100}+\dots+w_0^3a_{0003})\subseteq \mathbb{P}^{19}\)

Defines a morphism  $$\mathbb{P}^3\to \check{\mathbb{P}}^{19}$$ (the Veronese map of deg 3)

\textbf{Caution} Kleiman–Bertini theorem ensures these hyperplane are general enough!
          (for a general \(p\in \mathbb{P}^3\), \(\Pi_p\) and \(\Omega_F\) intersect transversally)

\subsubsection{Coordinates for 15 points}
\label{sec:org944217c}

\begin{minted}[breaklines,linenos,numbersep=3pt,mathescape=true,frame=lines,framesep=2mm]{julia}
@var P[1:15,1:4]
# parameters to do monodromy loops
p = P[:,4]
\end{minted}

\subsubsection{Equations with parameters}
\label{sec:orgac5be9d}

\begin{minted}[breaklines,linenos,numbersep=3pt,mathescape=true,frame=lines,framesep=2mm]{julia}
Eq2 = [subs(F, [x,y,z,w] => A*P[i,:]) for i in 1:15];
\end{minted}

\subsubsection{Computing initial pair solution-parameter}
\label{sec:orgebd9112}

\begin{enumerate}
\item Fix random parameters
\label{sec:org194e4d9}

We pick 15 random points (and leave one coordinate per point unspecified)

\begin{minted}[breaklines,linenos,numbersep=3pt,mathescape=true,frame=lines,framesep=2mm]{julia}
P₁₂₃ = randn(ComplexF64, 15, 3);
\end{minted}

\item System of eq. for monodromy
\label{sec:orgcfae20f}

\begin{minted}[breaklines,linenos,numbersep=3pt,mathescape=true,frame=lines,framesep=2mm]{julia}
Eq2_mono = subs(Eq2, A[4,4] => 1, vec(P[:,1:3]) => vec(P₁₂₃));
\end{minted}

\item Initial solution: random linear transformation.
\label{sec:org5a3a5cb}

\begin{minted}[breaklines,linenos,numbersep=3pt,mathescape=true,frame=lines,framesep=2mm]{julia}
A₀ = randn(ComplexF64, 4, 4)
A₀[4,4] = 1
# Initial solution for the affine equations.
a₀ = vec(A₀)[1:15]
\end{minted}

\item Compute initial parameters: the last coordinate
\label{sec:org41ec9ea}

\begin{minted}[breaklines,linenos,numbersep=3pt,mathescape=true,frame=lines,framesep=2mm]{julia}
p₀ = map(Eq2_mono) do f
    # substitute and compute root
    subs_f = subs(f, vec(A) => vec(A₀))
    last(roots(coefficients(subs_f, variables(subs_f)) |> reverse))
end; # initial value of the parameters
\end{minted}
\end{enumerate}

\subsubsection{Solving monodromy}
\label{sec:org3f552a9}

\begin{minted}[breaklines,linenos,numbersep=3pt,mathescape=true,frame=lines,framesep=2mm]{julia}
System2 = System(Eq2_mono, variables=vec(A)[1:15], parameters=p)
res = monodromy_solve(System2, a₀, p₀)
\end{minted}

\#\# Solutions found: 96120 	 Time: 0:45:35
\#\# MonodromyResult
\#\# \texttt{================================}
\#\# • 96120 solutions (0 real)
\#\# • return code → heuristic\textsubscript{stop}
\#\# • 1263490 tracked paths
\#\# • seed → 186832

\section{Now what?}
\label{sec:org51017ca}

Can we do something more than simply have faith in the computations?

Yes, we can! (what a time to be alive\ldots{})

\subsection{Alpha certificate (lower bound)}
\label{sec:org03a03a8}

\textbf{Approximate zero} Given a polynomial system \(F=0\),
a point \(p\in \mathbb{C}^n\) is an \emph{approximate zero} of \(F=0\)
if the Newton method starting at \(p\) converges quadraticly to a (unique) zero of \(F=0\).

Smale's \(\alpha\) theorem:

\begin{itemize}
\item Our case: We just want to ensure that our points actually approximate distinct zeros!
\end{itemize}


\subsection{Trace test (upper bound)}
\label{sec:org2769c00}

\href{trace-test.pdf}{Trace test}

\section{Classical proof}
\label{sec:org9c5aa7f}

Previous work on plane curves
\href{aluffi.pdf}{LINEAR ORBITS OF SMOOTH PLANE CURVES}

Using the same techniques
\href{2002-Article Text-6177-1-10-20200909.pdf}{Towards the degree of the PGL(4)-orbit of a cubic surface}

Classical proof with modern techniques
\href{A UNIVERSAL FORMULA FOR COUNTING CUBIC SURFACES.pdf}{A UNIVERSAL FORMULA FOR COUNTING CUBIC SURFACES}



























































\subsection{Misc}
\label{sec:orgec2e0a1}

Lorde
Jorja Smith
Erika de Casier
Japanese breakfast
Ela Minus
\end{document}
